% TEMPLATE-MILC-v3.tex - plantilla maestra UNICA de las guias MILC v3.
%
% La usan las cuatro familias de guias, cada una con su driver:
%   · Grados 6-11          -> scripts/build-guias-g11.py
%   · Bebras               -> scripts/build-guias-bebras.py
%   · Semillero            -> scripts/build-guias-semillero.py
%   · Territorio interior  -> scripts/build-guias-territorio-interior.py
%
% Antes habia tres copias casi identicas (~400 lineas iguales, 6 distintas):
% cada mejora habia que aplicarla tres veces, y en la practica no se hacia
% (el motor de recursos vivia solo en dos de ellas). Lo que varia entre
% familias esta parametrizado en 6 placeholders que cada driver llena
% (se nombran aqui SIN su sintaxis real: el builder reemplaza por texto plano,
% tambien dentro de los comentarios, y un valor multilinea rompe el preambulo):
%
%   HEADER_IZQ        encabezado izquierdo de pagina
%   HEADER_DER        encabezado derecho de pagina
%   PORTADA_ETIQUETA  rotulo sobre el numero grande (GRADO/MOMENTO/MODULO)
%   PORTADA_SERIE     linea de serie de la portada
%   PORTADA_ID        identificador de la guia en la portada
%   PORTADA_CIRCULO   el circulito de grado (°), o vacio si no aplica
%   PDF_TITULO        titulo en texto plano (sin \textbf ni comillas TeX) para
%                     los metadatos del PDF (/Title); lo produce lib_xelatex.texto_pdf
%
% Composicion (auditoria 2026-09-03): las bandas de estacion piden espacio
% (\Needspace*) y prohiben el salto justo despues (\nopagebreak), asi no quedan
% huerfanas al pie; las cajas largas (softbox, drawbox, infoband, puentebox) son
% `breakable`, asi una caja de media pagina ya no arrastra todo a la siguiente
% dejando la anterior al 40 %. Todo texto sobre mostaza o verde va en milcNegro
% (blanco sobre #E5B400 da 1,93:1 y sobre #58A923 2,95:1; ninguno pasa WCAG AA).
% El resto de claves las documenta content/guias/_SCHEMA.md. El script valida
% que no queden placeholders sin reemplazar antes de escribir el archivo final.

% Etiquetado PDF (latex-lab): /Lang, estructura y texto alternativo de las
% figuras, para lectores de pantalla. Debe ir ANTES de \documentclass.
\DocumentMetadata{lang=es-CO,tagging=on}
\documentclass[11pt,letterpaper]{article}
% headheight=24pt: la cabecera derecha lleva dos lineas (identidad de la guia
% y estacion actual). headsep se acorta en la misma medida para que el bloque
% de texto quede exactamente donde estaba con headheight=12pt.
\usepackage[margin=2.0cm,headheight=24pt,headsep=13pt]{geometry}
\usepackage[table]{xcolor}
\usepackage{fontspec}
\usepackage[spanish]{babel}
\usepackage{tikz}
% Librerías TikZ para los diagramas de las guías (bloque `recursos:`):
%  · babel  → babel en español vuelve activos caracteres como «>», y TikZ los usa
%             en sus opciones (>=stealth, ->). Sin esta librería, cualquier
%             diagrama con flechas revienta con «\language@active@arg> has an
%             extra }». Debe ir ANTES de que se dibuje nada.
%  · positioning        → sintaxis «below left=2pt and 3pt».
%  · decorations.pathreplacing → llaves ({brace}) para señalar rangos.
\usetikzlibrary{babel,positioning,decorations.pathreplacing}
\usepackage{graphicx}
% Circuitos: el trabajo de electrónica se hace hoy con micro:bit y sus sensores
% integrados (MakeCode, sin cableado), así que no se necesitan esquemáticos.
% Si algún día entran montajes con protoboard, basta descomentar esta línea:
% los diagramas .tex/.tikz declarados en `recursos:` ya se incrustan solos.
% \usepackage{circuitikz}
\usepackage[most]{tcolorbox}
\usepackage{tabularx}
\usepackage{array}
\usepackage{fancyhdr}
\usepackage{titlesec}
\usepackage[document]{ragged2e}  % bandera a la derecha en todo el cuerpo
\usepackage{enumitem}
\usepackage{microtype}
\usepackage{etoolbox}
\usepackage{needspace}
% hyperref al final del preambulo: metadatos del PDF (titulo, autor, idioma,
% familia), marcadores por estacion (\pdfbookmark) y enlaces sin recuadro.
\usepackage[hidelinks,
  pdftitle={Irse o quedarse: lo que pesa no es la distancia, es la red},
  pdfauthor={PhD. Álvaro Cárdenas Orozco},
  pdfsubject={Territorio interior · Educación socioemocional · Grado 11° · Momento 5 · MILC v3},
  pdflang={es-CO},
  bookmarksopen=true,bookmarksnumbered=false]{hyperref}

% Helvetica Neue no trae la flecha U+2192, asi que cada «→» escrito literalmente
% en el YAML se compilaba a NADA, en silencio (462 flechas en 61 guias: rutas del
% tipo «paleta → tipografia → lockup» salian rotas). Mapeamos el caracter a su
% version matematica, que si tiene glifo. Asi el autor puede seguir escribiendo
% «→» comodamente en el YAML.
\usepackage{newunicodechar}
\newunicodechar{→}{\ensuremath{\rightarrow}}
% Mismo problema con los simbolos de marca y vinetas: el «✓» de «una palomita ✓»
% salia como cuadrito vacio en el PDF de todas las guias que lo usaban.
\usepackage{pifont}
\newunicodechar{✓}{\ding{51}}
\newunicodechar{✗}{\ding{55}}
\newunicodechar{✔}{\ding{52}}
\newunicodechar{•}{\textbullet}
\newunicodechar{≥}{\ensuremath{\geq}}
\newunicodechar{≤}{\ensuremath{\leq}}

\setmainfont{Helvetica Neue}
\setsansfont{Helvetica Neue}
\setlength{\parindent}{0pt}
\setlength{\parskip}{6pt}
% Sin particion de palabras: en 6.º-7.º una palabra cortada por guion al final
% de linea («Comuni-cacion») frena la lectura mucho mas que una linea corta.
\hyphenpenalty=10000
\exhyphenpenalty=10000
\pretolerance=2000
\tolerance=3000
\setlength{\emergencystretch}{3em}
\linespread{1.10}
\renewcommand{\arraystretch}{1.25}
\setlist{leftmargin=5mm,itemsep=1mm,topsep=1mm}
\newcolumntype{Y}{>{\RaggedRight\arraybackslash}X}

\definecolor{milcMagenta}{HTML}{E6007E}
\definecolor{milcVino}{HTML}{5A0038}
\definecolor{milcTurquesa}{HTML}{0093A5}
\definecolor{milcVerde}{HTML}{58A923}
\definecolor{milcMostaza}{HTML}{E5B400}
\definecolor{milcOcre}{HTML}{B86600}
\definecolor{milcNegro}{HTML}{241816}
\definecolor{milcGris}{HTML}{F7F7F4}
\definecolor{milcLinea}{HTML}{E8E2DD}
\definecolor{milcGradePrimary}{HTML}{0D9488}
\definecolor{milcGradeDark}{HTML}{115E59}
\definecolor{milcGradeSoft}{HTML}{CCFBF1}
\definecolor{milcGradeAccent}{HTML}{CFE7FF}
\definecolor{milcGradeText}{HTML}{FFFFFF}
\color{milcNegro}

\pagestyle{fancy}
\fancyhf{}
% Cabecera de dos lineas: identidad de la guia arriba y, a la derecha y en
% gris, la estacion en curso (\markright desde \milcestacion). Antes de la
% primera estacion la segunda linea va vacia; el \strut mantiene la altura
% para que la linea de arriba no baile entre paginas.
\lhead{\small Territorio interior · Educación socioemocional\\[-1pt]{\footnotesize\strut}}
\rhead{\small Grado 11° · Momento 5 · MILC v3\\[-1pt]{\footnotesize\strut\color{milcNegro!65}\rightmark}}
\cfoot{\small\thepage}
\renewcommand{\headrulewidth}{0pt}

% Sobre mostaza (#E5B400) y verde (#58A923) el texto va en milcNegro; sobre el
% resto de la paleta, en blanco. Se usa dentro de fonttitle/fontupper, donde un
% \color posterior manda sobre coltitle.
\newcommand{\milctintasobre}[1]{%
  \ifboolexpr{test{\ifstrequal{#1}{milcMostaza}} or test{\ifstrequal{#1}{milcVerde}}}%
    {\color{milcNegro}}{\color{white}}}

\newtcolorbox{titlebox}[1]{
  enhanced,colback=#1,colframe=#1,arc=7mm,boxrule=0pt,
  left=7mm,right=7mm,top=3mm,bottom=3mm,
  before skip=8pt,after skip=8pt,
  fontupper=\sffamily\bfseries\Large\color{white}
}

\newtcolorbox{softbox}[3]{
  enhanced,breakable,pad at break=3mm,
  colback=#2,colframe=#1,borderline west={4pt}{0pt}{#1},
  arc=2mm,boxrule=.4pt,left=4mm,right=4mm,top=3mm,bottom=3mm,
  before skip=6pt,after skip=8pt,title={#3},coltitle=white,
  colbacktitle=#1,fonttitle=\sffamily\bfseries\milctintasobre{#1},fontupper=\RaggedRight,
  attach boxed title to top left={xshift=3mm,yshift=-2mm},
  boxed title style={arc=2mm, boxrule=0pt}
}

\newtcolorbox{drawbox}[2]{
  enhanced,breakable,pad at break=4mm,
  colback=white,colframe=#1,arc=4mm,boxrule=1pt,
  left=5mm,right=5mm,top=4mm,bottom=4mm,before skip=8pt,after skip=8pt,
  title={#2},coltitle=white,colbacktitle=#1,fonttitle=\sffamily\bfseries\milctintasobre{#1},
  fontupper=\RaggedRight,
  attach boxed title to top left={xshift=4mm,yshift=-2mm},
  boxed title style={arc=3mm, boxrule=0pt}
}

\newtcolorbox{infoband}[1]{
  enhanced,breakable,pad at break=2mm,colback=#1!8,colframe=#1!50,arc=2mm,boxrule=.5pt,
  left=4mm,right=4mm,top=2mm,bottom=2mm,before skip=4pt,after skip=4pt,
  fontupper=\small\RaggedRight
}

% Puente narrativo entre secciones (contrato editorial Conéctate).
% Da continuidad: "Acabas de X. Ahora vas a Y porque Z."
% Visualmente: banda izquierda en color del grado, texto en itálica pequeña.
\newtcolorbox{puentebox}{
  enhanced,breakable,pad at break=2mm,colback=milcGradeSoft,colframe=milcGradeDark,
  borderline west={3pt}{0pt}{milcGradeDark},
  arc=0mm,boxrule=0pt,left=5mm,right=4mm,top=2.5mm,bottom=2.5mm,
  before skip=4pt,after skip=8pt,
  fontupper=\itshape\small\color{milcNegro}\RaggedRight
}
% La flecha va en modo matematico a proposito: Helvetica Neue NO tiene el glifo
% U+2192, asi que \textrightarrow se compilaba a NADA («Missing character: There
% is no → in font Helvetica Neue») y el puente salia sin flecha. En modo
% matematico la flecha viene de la fuente matematica y si se dibuja.
\newcommand{\puente}[1]{%
  \begin{puentebox}%
    \textcolor{milcGradeDark}{$\rightarrow$}\hspace{2mm}#1%
  \end{puentebox}%
}

\newcommand{\linead}[1]{\noindent\color{milcLinea}\rule{#1}{0.5pt}\color{milcNegro}}
\newcommand{\respuesta}[1]{\par\vspace{2mm}\foreach \n in {1,...,#1}{\noindent\color{milcLinea}\rule{\linewidth}{0.5pt}\par\vspace{5mm}}\color{milcNegro}}
\newcommand{\checkbox}{\textcolor{milcGradeDark}{\fbox{\phantom{x}}}\hspace{5pt}}

% ─── Listas aireadas para las guías socioemocionales ────────────────────
% Componer las verificaciones como «\checkbox texto\\» deja el texto que salta
% de línea debajo del cuadrito y apelmaza el bloque. Estos entornos alinean la
% continuación y airean los ítems. Los usa Territorio Interior; cualquier
% familia puede hacerlo.
\newlist{checklista}{itemize}{1}
\setlist[checklista]{
  label=\textcolor{milcGradeDark}{\fbox{\phantom{x}}},
  leftmargin=9mm, labelsep=3.2mm, itemsep=2.4mm,
  topsep=2.5mm, parsep=0pt, partopsep=0pt, before=\RaggedRight
}
\newlist{pasoslista}{enumerate}{1}
\setlist[pasoslista]{
  label=\textcolor{milcGradeDark}{\sffamily\bfseries\arabic*.},
  leftmargin=8mm, labelsep=3mm, itemsep=2.8mm,
  topsep=1mm, parsep=0pt, partopsep=0pt, before=\RaggedRight
}

\newcommand{\phasecard}[4]{
\begin{tcolorbox}[enhanced,colback=#3,colframe=#2,arc=3mm,boxrule=.5pt,height=3.05cm,valign=center,halign=center,left=1.5mm,right=1.5mm,top=2mm,bottom=2mm]
{\sffamily\bfseries\fontsize{22}{24}\selectfont\textcolor{#2}{#1}}\par
{\sffamily\bfseries\small #4}
\end{tcolorbox}
}

% ── Piezas del rediseno 2026 (legibilidad 6.º-11.º) ───────────────────
% Banda de estacion: reemplaza el titlebox macizo. Titulo a la izquierda,
% dato practico (tiempo/modalidad) a la derecha, en una sola linea.
% Nunca huerfana: pide 9 lineas libres (o salta de pagina rellenando con
% \vfill) y prohibe el salto justo despues de la banda. Nueve y no seis:
% la caja partible que sigue necesita ~80pt (titulo adosado, relleno y dos
% lineas) para arrancar en la misma pagina; con menos, tcolorbox la manda
% entera a la siguiente y la banda se queda sola. El penalty 10000 va
% en `after`, ANTES del espacio posterior: un `after skip` (aunque sea 0pt)
% es un \vskip tras la caja, y ese glue es un punto de corte legal que un
% \nopagebreak escrito despues de \end{tcolorbox} ya no cierra. Cada banda
% deja marcador PDF y actualiza la cabecera derecha.
\newcounter{milcestacion}
\newcommand{\milcestacion}[2]{%
\Needspace*{9\baselineskip}%
\stepcounter{milcestacion}%
\pdfbookmark[1]{#2}{milcestacion\themilcestacion}%
\markright{#2}%
\begin{tcolorbox}[enhanced,colback=#1,colframe=#1,arc=2mm,boxrule=0pt,
  left=5mm,right=5mm,top=2.6mm,bottom=2.6mm,before skip=11pt,
  after={\par\nopagebreak\vskip7pt}]
{\sffamily\bfseries\fontsize{15}{18}\selectfont\milctintasobre{#1}#2}
\end{tcolorbox}}

% Tarjeta de paso numerada: sustituye las tablas «Pilar / Aplicacion», que
% eran formato de consulta docente, no de estudiante. El numero va en un
% circulo sobre el borde para que la secuencia se lea de un vistazo.
\newcommand{\milcpaso}[3]{%
\Needspace{4\baselineskip}%
\begin{tcolorbox}[enhanced,colback=white,colframe=milcLinea,arc=1.5mm,boxrule=.9pt,
  left=14mm,right=4mm,top=3mm,bottom=3mm,before skip=4pt,after skip=4pt,
  fontupper=\RaggedRight,
  overlay={\node[anchor=north west,circle,fill=#2,text=white,inner sep=0pt,
    minimum size=7.4mm,font=\sffamily\bfseries\small]
    at ([xshift=3.6mm,yshift=-3.2mm]frame.north west) {#1};}]
#3
\end{tcolorbox}}

% Casilla marcable de tamano util con lapiz (la anterior era \fbox{\phantom{x}},
% de alto variable segun el texto que la seguia).
\newcommand{\milcmarca}{\framebox[3.8mm]{\rule{0pt}{3.4mm}}}

% Fila marcable de lista suelta (checks de la estacion 1).
\newcommand{\milccheckrow}[1]{%
\par\vspace{1.8mm}\noindent
\makebox[6.4mm][l]{\raisebox{-.55mm}{\textcolor{milcNegro}{\milcmarca}}}%
\parbox[t]{\dimexpr\linewidth-6.4mm\relax}{\RaggedRight #1}\par}

% Celda marcable dentro de la rubrica (tabla de autoevaluacion). Misma
% sangria francesa, pero pensada para vivir dentro de una columna X.
\newcommand{\milcopcion}[1]{%
\makebox[5.2mm][l]{\raisebox{-.55mm}{\textcolor{milcNegro}{\milcmarca}}}%
\parbox[t]{\dimexpr\linewidth-5.2mm\relax}{\RaggedRight #1}}

% ── Recursos visuales de la guía ──────────────────────────────────────
% Declarados en el bloque `recursos:` del YAML y emitidos por el builder.
% \guiaFigura   → imagen rasterizada o PDF (foto, carta celeste, montaje).
% \guiaDiagrama → fuente TikZ/circuitikz (.tex) incrustada como vector:
%                 es la vía para circuitos y esquemáticos editables.
% [#1] = texto alternativo (alt) · #2 = ruta relativa al .tex · #3 = pie
% El alt llega al PDF etiquetado (/Alt) y lo lee el lector de pantalla.
\newcommand{\guiaFigura}[3][]{%
\begin{center}
\includegraphics[width=0.88\linewidth,height=0.38\textheight,keepaspectratio,alt={#1}]{#2}\\[1.5mm]
{\footnotesize\itshape\textcolor{milcNegro!75}{#3}}
\end{center}
}

\newcommand{\guiaDiagrama}[2]{%
\begin{center}
\input{#1}\\[1.5mm]
{\footnotesize\itshape\textcolor{milcNegro!75}{#2}}
\end{center}
}

\begin{document}
\thispagestyle{empty}

% ── Portada ───────────────────────────────────────────────────────────
% Antes: un rectangulo de color a pagina completa. Se veia bien en pantalla,
% pero en la fotocopiadora del colegio se comia el toner y salia gris sucio.
% Ahora la banda ocupa el tercio superior y el resto va en blanco.
\begin{tikzpicture}[remember picture,overlay]
  \path[fill=milcGradePrimary]
    (current page.north west) rectangle ([yshift=-10.6cm]current page.north east);

  \node[anchor=north west,text=milcGradeText] at ([xshift=2.0cm,yshift=-1.55cm]current page.north west)
    {\sffamily\bfseries\fontsize{12}{14}\selectfont MOMENTO};
  \node[anchor=north west,text=milcGradeText,opacity=.85] at ([xshift=2.0cm,yshift=-2.35cm]current page.north west)
    {\sffamily\bfseries\fontsize{8.5}{10}\selectfont TERRITORIO INTERIOR · SOCIOEMOCIONAL};
  \node[anchor=north east,text=milcGradeText,opacity=.85,align=right] at ([xshift=-2.0cm,yshift=-1.55cm]current page.north east)
    {\sffamily\bfseries\fontsize{9}{12}\selectfont I.E. Sor María Juliana\\Cartago, Valle del Cauca};

  \node[anchor=west,text=milcGradeText] (gradonum) at ([xshift=1.85cm,yshift=-7.1cm]current page.north west)
    {\sffamily\bfseries\fontsize{112}{112}\selectfont 5};
  % El circulito de grado (°) solo aplica a las guias de grado («11°»); en
  % Bebras y Semillero el numero grande es un momento/modulo, no un grado.
  % Se posiciona relativo al nodo `gradonum`, no en coordenadas absolutas.
  

  \node[anchor=west,text=milcGradeText,text width=11.4cm,align=left] at ([xshift=1.5cm,yshift=0.5cm]gradonum.east)
    {
     {\sffamily\bfseries\fontsize{9}{11}\selectfont\MakeUppercase{Territorio interior · Socioemocional}}\\[3mm]
     {\sffamily\bfseries\fontsize{21}{25}\selectfont\setlength{\baselineskip}{27pt}Irse\\[4pt]o quedarse\\[4pt]lo que pesa es la red}\\[3.5mm]
     {\sffamily\bfseries\fontsize{15}{18}\selectfont Grado 11° · Momento 5}
    };
\end{tikzpicture}

\vspace*{9.4cm}
\vfill

{\sffamily\bfseries\fontsize{8}{10}\selectfont\textcolor{milcGradeDark}{LO QUE TE LLEVAS HOY}}

\begin{tcolorbox}[enhanced,colback=white,colframe=milcNegro,arc=2mm,boxrule=1.6pt,
  left=5mm,right=5mm,top=4mm,bottom=4mm,before skip=3pt,after skip=12pt,
  fontupper=\RaggedRight\fontsize{11}{15.5}\selectfont]
Tu balance del irse: las dos columnas de lo que ganas y lo que pierdes en cada opción, el inventario de la red que tienes aquí y de la que tendrías allá, la prueba del arrepentimiento a diez años, y la conversación con alguien que se fue.
\end{tcolorbox}

\begin{tcolorbox}[enhanced,colback=milcGris,colframe=milcGris,arc=2mm,boxrule=0pt,
  left=5mm,right=5mm,top=3.5mm,bottom=3.5mm,after skip=12pt]
{\sffamily\bfseries\fontsize{8}{10}\selectfont\textcolor{milcNegro!55}{ESTA GUÍA ES DE}}\\[3mm]
\begin{tabularx}{\linewidth}{X p{3.2cm} p{3.2cm}}
\linead{\linewidth} & \linead{3.0cm} & \linead{3.0cm}\\[-1mm]
{\fontsize{8.5}{10}\selectfont\textcolor{milcNegro!55}{Nombre completo}} &
{\fontsize{8.5}{10}\selectfont\textcolor{milcNegro!55}{Grupo}} &
{\fontsize{8.5}{10}\selectfont\textcolor{milcNegro!55}{Fecha}}\\
\end{tabularx}
\end{tcolorbox}

\vfill

\noindent\textcolor{milcLinea}{\rule{\linewidth}{1pt}}\\[2mm]
{\sffamily\bfseries\fontsize{9.5}{12}\selectfont Autor: Dr. Álvaro Cárdenas Orozco}\\[1mm]
{\sffamily\fontsize{8.5}{11}\selectfont\textcolor{milcNegro!55}{Metodología MILC · Apertura ancestral · Migración y querencia · Triángulo Dussel-Estoicismo-Floridi}}

\newpage

\section*{Apertura: Irse o quedarse: lo que pesa no es la distancia, es la red}

\begin{softbox}{milcGradeDark}{milcGradeSoft}{Saber ancestral}
Colombia es un país que se ha movido mucho por dentro. \emph{Desde 1930 la ciudad se volvió mucho más atractiva y empezó a recibir migrantes del campo; hacia 1947 las ciudades comenzaron además a recibir gente desplazada por la violencia,} \textbf{y se formaron barrios enteros de recién llegados.} \textbf{Y hay un mecanismo, documentado en los estudios sobre migración, que explica quién sobrevivía bien a esa llegada:} \emph{las \textbf{redes de paisanos}}. \textbf{Las redes se iban configurando y fusionando para \emph{estimular, facilitar y apoyar} la llegada de quienes venían, y también su integración social,} \emph{sobre todo cuando eran recién llegados.} \textbf{El dato que importa es este:} \emph{la adaptación de quienes llegaron \textbf{después} de los pioneros fue \textbf{más fácil}, porque en su condición de recién llegados \textbf{encontraban paisanos ya establecidos} que les brindaban apoyo o estaban dispuestos a darles información.} \emph{Fíjate en lo que eso dice y en lo que no dice.} \textbf{No dice que unos migrantes fueran más capaces que otros.} \emph{Dice que la diferencia estaba \textbf{afuera de la persona}: en si al llegar había alguien.} \textbf{Los pioneros la pasaron mucho peor no por ser distintos, sino por llegar primero.}
\end{softbox}

\begin{softbox}{milcTurquesa}{milcGris}{Saber contemporáneo}
¿Y de qué se arrepiente la gente al final: de lo que hizo o de lo que no hizo? \textbf{Depende de cuándo se le pregunte, y el patrón es muy consistente.} \textbf{Thomas Gilovich} y \textbf{Victoria Medvec} describieron lo que llamaron \textbf{el patrón temporal del arrepentimiento}: \emph{a \textbf{corto plazo} la gente se arrepiente más de \textbf{lo que hizo}; a \textbf{largo plazo}, más de \textbf{lo que no hizo}} (Gilovich y Medvec, 1995). \textbf{Y explicaron por qué:} \emph{el escozor de una acción de la que uno se arrepiente \textbf{disminuye relativamente rápido}} ---entre otras cosas porque uno hace «trabajo de reparación»: arregla, compensa, le encuentra sentido---, \emph{mientras que \textbf{el dolor de lo que no se hizo permanece}}, porque no hay nada que reparar ni desenlace que asimilar. \textbf{Encontraron incluso que las emociones son distintas:} \emph{lo que uno hizo produce emociones «calientes» ---rabia, vergüenza---; lo que uno no hizo produce \textbf{nostalgia y desconsuelo}}. \emph{Es decir:} \textbf{equivocarse haciendo duele más ahora; no haberlo intentado duele más durante años.}
\end{softbox}

\begin{softbox}{milcMostaza}{milcGris}{Pregunta-puente}
\textit{Los que llegaron segundos la pasaron mejor \textbf{porque había alguien esperándolos}. \textbf{Si te fueras el otro año, ¿quién estaría allá} \ldots \textbf{y quién se quedaría aquí}?}
\end{softbox}

\begin{softbox}{milcVerde}{milcGris}{Saber y saber hacer}
Al terminar podrás: (1) \textbf{entender} que lo que decide una migración no es la distancia ni el mérito personal, sino la red que uno encuentra o construye al llegar; (2) \textbf{explicar} la asimetría del arrepentimiento y usarla para mirar tu propia decisión a largo plazo; (3) \textbf{inventariar} la red que tienes aquí y la que tendrías allá, con nombres; (4) \textbf{hacer} un balance honesto de las dos opciones, sabiendo que quedarse también es una decisión que se puede defender.
\end{softbox}



\small
\begin{tabularx}{\textwidth}{>{\bfseries}p{3.0cm}Y}
\rowcolor{milcGradePrimary!12}Autor & Dr. Álvaro Cárdenas Orozco\\
DBA & Comprende los fenómenos que afectan a su territorio, regula sus emociones ante ellos y acompaña a otros con empatía (transversal 6°--11°, articulado con la política «Escuela, territorio de vida» del MEN, Resolución 006519 de 2025, y la Ley 1620 de 2013)\\
Referentes & Primera ayuda psicológica (OMS, 2011/2012) · Directrices IASC (2007) · USGS y Servicio Geológico Colombiano · Pueblo nasa y la minga (CRIC) · Dussel · Estoicismo · Floridi\\
Producto final & Tu balance del irse: las dos columnas de lo que ganas y lo que pierdes en cada opción, el inventario de la red que tienes aquí y de la que tendrías allá, la prueba del arrepentimiento a diez años, y la conversación con alguien que se fue.\\
\end{tabularx}
\normalsize

\puente{En el momento 4 miraste el oficio. \textbf{Hoy miras el lugar}, \emph{que es la otra mitad de la decisión que tienes encima ---y la que casi nunca se piensa con método---}.}

\milcestacion{milcGradeDark}{Ruta MILC: así vamos a aprender}
\begin{tabularx}{\textwidth}{>{\centering\arraybackslash}X>{\centering\arraybackslash}X>{\centering\arraybackslash}X>{\centering\arraybackslash}X}
\phasecard{1}{milcVerde}{milcGris}{Escucha\\[-1mm]\small Observo, escucho y pregunto.} &
\phasecard{2}{milcTurquesa}{milcGris}{Sistematización\\[-1mm]\small Miro la red, no la distancia.} &
\phasecard{3}{milcMagenta}{milcGris}{Praxis\\[-1mm]\small Construyo y aplico.} &
\phasecard{4}{milcVino}{milcGris}{Evaluación\\[-1mm]\small Reflexión liberadora.}
\end{tabularx}


\puente{Primero, sin idealizar ninguna de las dos: qué te empuja a irte y qué te ata aquí.}

\milcestacion{milcVerde}{Estación 1 de 3 · Escucha}

\begin{softbox}{milcVerde}{milcGris}{Escena para escuchar}
\textbf{Actividad 1 · IDENTIFICA --- Lo que empuja y lo que ata.} \textbf{Nadie va a leer esta página.} \textbf{Primera parte.} Escribe \textbf{qué te empujaría a irte} de donde vives ---estudiar algo que no hay aquí, trabajo, ganas de conocer, salir de una situación, que aquí no pasa nada--- \textbf{y qué te ata} ---la familia, alguien en particular, la plata, el miedo, que aquí sabes moverte---. \textbf{Sin adornar ninguna de las dos listas.} \textbf{Segunda parte.} De la gente que conoces \textbf{que se fue}: \emph{¿a quiénes les fue bien y a quiénes no?} \textbf{Y la pregunta que importa:} \emph{¿había alguien esperándolos allá?} \textbf{Tercera parte.} Haz dos listas de \textbf{nombres}: \emph{a quién tienes aquí} ---la gente con la que cuentas de verdad--- y \emph{a quién tendrías allá}. \textbf{Cuenta cuántos hay en cada una.} \textbf{Y la última:} \emph{si te fueras, ¿quién te haría más falta, y quién sentiría más tu falta?} \emph{No siempre es la misma persona.}
\end{softbox}

{\sffamily\bfseries\fontsize{8}{10}\selectfont\textcolor{milcNegro!55}{MARCA CUANDO LO HAYAS HECHO}}
\milccheckrow{Escribiste las dos listas ---lo que empuja y lo que ata--- sin adornarlas.}
\milccheckrow{Anotaste si había alguien esperando a quienes se fueron.}
\milccheckrow{Contaste \textbf{nombres} en la red de aquí y en la de allá.}

\begin{infoband}{milcVerde}


\textbf{En tu cuaderno (Actividad 1 · IDENTIFICA).} Título: «Actividad 1. Lo que empuja y lo que ata». \textbf{Formato:} las dos listas + los casos que conoces con la pregunta de la red + los nombres de aquí y de allá con sus números + a quién extrañarías. \textbf{Extensión:} media página. \textbf{Sabes que terminaste cuando} tienes \textbf{los dos números}. Esta página es tuya: \textbf{nadie más tiene que leerla si tú no quieres}.
\end{infoband}

\puente{Ya lo tienes. Ahora, por qué lo que pesa es la red, y qué dice la investigación sobre de qué se arrepiente uno con los años.}

\milcestacion{milcTurquesa}{Fase 2 · Sistematización: entender la partida}

\begin{softbox}{milcTurquesa}{milcGris}{Cuatro claves sobre irse}
Cuatro claves para decidir una partida con algo más que ganas o miedo.
\end{softbox}

\milcpaso{1}{milcTurquesa}{\textbf{Lo primero que se pierde no es el lugar: es la gente que sabía quién eras.} \emph{Cuando alguien se va, se imagina extrañando el paisaje, la comida, el clima.} \textbf{Y lo que de verdad pesa es otra cosa:} \emph{en el lugar de donde uno viene, la gente ya sabe cómo es uno} ---no hay que explicarse, no hay que empezar de cero, alguien se acuerda de cosas que a uno le pasaron---. \textbf{Allá, uno es un desconocido}, \emph{y serlo cansa mucho más de lo que parece.} \emph{Por eso los primeros meses son duros incluso cuando todo lo demás sale bien:} \textbf{no es que la ciudad sea hostil ---es que nadie tiene historia con uno---.} \emph{Y eso también dice qué hay que construir allá, y que se construye lento.}}
\milcpaso{2}{milcTurquesa}{\textbf{El que llega segundo llega mejor, y no es por mérito.} \emph{Los estudios sobre migración en Colombia lo muestran con claridad:} \textbf{la adaptación de quienes llegaban después de los pioneros era más fácil, porque encontraban paisanos ya establecidos} que les daban apoyo o información. \emph{Léelo al derecho y al revés.} \textbf{Al derecho:} \emph{si vas a irte, la pregunta más útil no es «¿a qué universidad?» sino \textbf{«¿quién está allá?»}} ---un primo, un paisano, un egresado del colegio, alguien de la iglesia, un compañero del año pasado---. \textbf{Al revés, y esto también te toca:} \emph{si alguien de tu pueblo llega a donde tú estés, \textbf{tú eres el que ya está establecido}}, \textbf{y lo que hagas define cómo le va.} \emph{La red no aparece: alguien la sostiene.}}
\milcpaso{3}{milcTurquesa}{\textbf{La asimetría del arrepentimiento, y cómo usarla bien.} \emph{A corto plazo la gente se arrepiente más de lo que hizo; a largo plazo, de lo que no hizo} (Gilovich y Medvec, 1995). \textbf{La razón es que lo hecho se puede reparar y con el tiempo pierde fuerza, mientras que lo no hecho queda sin desenlace} ---y produce nostalgia y desconsuelo en vez de rabia---. \emph{Ahora la advertencia, porque esta idea se usa mal todo el tiempo:} \textbf{esto \emph{no} significa «lánzate a lo que sea».} \emph{Significa una cosa más precisa:} \textbf{cuando estés comparando dos malestares, acuérdate de que \emph{no estás pesando bien el segundo}} ---el dolor de no haberlo intentado no se siente hoy, y por eso no entra en la cuenta---. \emph{Es la misma corrección que hiciste en el momento 3 con el sesgo de impacto, aplicada al otro lado de la balanza.}}
\milcpaso{4}{milcTurquesa}{\textbf{Quedarse también es decidir, y se puede defender.} \emph{Circula la idea de que irse es progresar y quedarse es conformarse.} \textbf{Es falsa, y hace daño en un territorio como este:} \emph{si todos los que pueden se van, el lugar se queda sin quienes lo sostienen ---y eso es exactamente lo que ha pasado con muchos pueblos---.} \textbf{Quedarse por decisión y quedarse por no haber decidido son cosas distintas} ---esa es la lección del momento 3---, \emph{y la primera se puede defender con argumentos: la red que ya tienes, lo que puedes construir aquí, lo que le debes al lugar ---que es el momento 8---.} \textbf{Y hay una tercera opción que casi nadie considera a los diecisiete:} \emph{irse un tiempo y volver}. \textbf{Mucha gente valiosa de esta región hizo exactamente eso.}}

\begin{drawbox}{milcTurquesa}{El balance, en cinco preguntas}
\begin{pasoslista}
\item \textbf{¿Quién está allá?} \emph{Con nombres. Si la respuesta es «nadie», eso no lo impide ---lo encarece---.}
\item \textbf{¿Qué red tendría que construir, y en cuánto tiempo?} \emph{Cuenta meses, no semanas.}
\item \textbf{¿Es reversible?} \emph{Casi siempre sí, y saberlo baja la presión de acertar.}
\item \textbf{A los veintisiete, ¿de qué me arrepentiría más?} \emph{Acordándote de que lo no hecho pesa más con los años.}
\item \textbf{Si me quedo, ¿es una decisión o es lo que pasó?} \emph{Si es decisión, ¿cuál es el plan?}
\end{pasoslista}
\end{drawbox}

\begin{softbox}{milcMostaza}{milcGris}{Errores comunes y su corrección}
\textbf{Creer que irse es siempre progresar:} y que quedarse es conformarse. \textbf{Mirar solo lo académico y lo económico:} lo que decide es la red. \textbf{«Yo me las arreglo solo»:} los pioneros la pasaron peor, y no por ser peores. \textbf{Usar la asimetría para lanzarse a cualquier cosa:} sirve para pesar bien, no para no pensar. \textbf{Idealizar el lugar de destino:} nadie tiene historia contigo allá. \textbf{Idealizar el de origen si te vas:} la nostalgia edita. \textbf{No considerar irse y volver:} es una opción real y muy frecuente.
\end{softbox}

\begin{infoband}{milcTurquesa}


\textbf{En tu cuaderno (Actividad 2 · EXPLICA).} Título: «Actividad 2. El patrón del arrepentimiento». \textbf{Formato:} explica con tus palabras por qué lo hecho duele más ahora y lo no hecho más después, y \textbf{aplícalo} a tu propia decisión sin usarlo como excusa para no pensar. \textbf{Extensión:} media página. \textbf{Sabes que terminaste cuando} puedes explicar por qué \textbf{el que llega segundo la pasa mejor}.
\end{infoband}

\puente{Ya sabes qué mirar. Es hora del balance con nombres propios, y de hablar con alguien que ya se fue.}

\milcestacion{milcMagenta}{Fase 3 · Praxis: hacer el balance}

\begin{softbox}{milcMagenta}{milcGris}{Cómo se decide una partida}
Un balance con nombres propios, y una conversación con alguien que ya lo hizo. Eso es todo, y es mucho más de lo que hace casi todo el mundo.
\end{softbox}

\milcpaso{1}{milcMagenta}{\textbf{Las dos columnas (8 min).} Para \textbf{cada} opción ---irse y quedarse--- escribe \emph{qué ganas} y \emph{qué pierdes}. \textbf{Cuatro celdas, no dos.} \emph{Casi todo el mundo llena solo dos ---lo que gana yéndose y lo que pierde quedándose--- y así la balanza sale amañada.} \textbf{Obligarte a escribir lo que ganas quedándote y lo que pierdes yéndote es la mitad del ejercicio.}}
\milcpaso{2}{milcMagenta}{\textbf{El inventario de la red (8 min).} Los nombres de aquí y los de allá, y para los de allá anota \textbf{qué tan cerca son de verdad} ---no «un primo lejano», sino si le podrías escribir mañana---. \textbf{Y calcula:} \emph{si me voy, ¿qué red tendría que construir y en cuántos meses?} \emph{Piensa en dónde se conoce gente:} \textbf{clase, trabajo, deporte, iglesia, un grupo de estudio, el paisanaje.}}
\milcpaso{3}{milcMagenta}{\textbf{La prueba de los diez años (5 min).} \emph{A los veintisiete: ¿de qué te arrepentirías más?} \textbf{Escríbelo acordándote de la asimetría} ---lo no hecho pesa más con los años, y hoy no se siente---. \emph{Y sé honesto si la respuesta es «de haberme ido»: también pasa, y también es información.}}
\milcpaso{4}{milcMagenta}{\textbf{La conversación con quien se fue (fuera de clase).} \textbf{Habla con alguien que se haya ido de aquí} ---un egresado, un familiar, alguien del barrio--- \textbf{y hazle cuatro preguntas:} \emph{«¿quién había allá cuando llegaste?», «¿cuánto tardaste en sentirte menos solo?», «¿qué es lo que más extrañas?», «¿lo volverías a hacer?»}. \emph{Y si conoces a alguien que se quedó pudiendo irse, hazle las mismas al revés.} \textbf{Anota las dos respuestas:} \emph{este es el único momento del programa donde vas a tener las dos versiones de la misma decisión.}}

\begin{drawbox}{milcMagenta}{Antes de cerrar tu balance}
\begin{checklista}
\item Llenaste \textbf{las cuatro celdas}, incluidas las incómodas.
\item Tu red de allá tiene \textbf{nombres} y una marca de qué tan cerca son.
\item Calculaste \textbf{en cuántos meses} construirías una red nueva.
\item Respondiste la prueba de los diez años, aunque la respuesta te sorprenda.
\item Hablaste con \textbf{alguien que se fue} ---y ojalá con alguien que se quedó---.
\item Consideraste la opción de \textbf{irse y volver}.
\end{checklista}
\end{drawbox}

\begin{softbox}{milcMagenta}{milcGris}{Plantilla de trabajo}
\textbf{Las cuatro celdas.} \emph{Irme: gano \_\_\_\_ / pierdo \_\_\_\_. Quedarme: gano \_\_\_\_ / pierdo \_\_\_\_.} \\[1mm]
\textbf{La red.} \emph{«Allá están: \_\_\_\_. Tendría que construir \_\_\_\_ en \_\_\_\_ meses.»} \\[1mm]
\textbf{Para quien se fue.} \emph{«¿Quién había allá cuando llegaste? · ¿Cuánto tardaste en sentirte menos solo? · ¿Qué extrañas? · ¿Lo volverías a hacer?»}
\end{softbox}

\begin{infoband}{milcMagenta}


\textbf{En tu cuaderno (Actividad 3 · CREA).} Título: «Actividad 3. Mi balance del irse». \textbf{Formato:} las cuatro celdas + el inventario de las dos redes con nombres + la prueba de los diez años + las respuestas de la conversación. \textbf{Extensión:} una página. \textbf{Sabes que terminaste cuando} tienes \textbf{las cuatro celdas llenas} y \textbf{una conversación hecha}.
\end{infoband}

\puente{El producto tiene nombres adentro. \emph{Una red no se cuenta en abstracto: se cuenta en personas.}}

\milcestacion{milcMostaza}{Producto de la sesión}
\textbf{Balance del irse: las cuatro celdas, el inventario de la red con nombres y la conversación con quien ya se fue}.
\begin{softbox}{milcMostaza}{milcGris}{Criterios de calidad}
(1) Las cuatro celdas están llenas, incluidas las que suelen omitirse. (2) El inventario de red usa nombres propios y evalúa la cercanía real de cada uno. (3) Se estima el tiempo de construcción de una red nueva y dónde se conocería gente. (4) La prueba de los diez años está respondida con honestidad, sin usarla como consigna. (5) La conversación se hizo con alguien que efectivamente migró, y se registran sus respuestas.
\end{softbox}

\puente{Antes de cerrar, revisa lo que suele faltar: si tu balance solo tiene plata y universidades, dejaste por fuera lo que de verdad decide.}

\milcestacion{milcVino}{Evaluación liberadora · cinco dimensiones}

\begin{softbox}{milcVino}{milcGris}{Sentido de la evaluación}
Evaluar no es solo revisar una nota. Es reconocer qué aprendí, qué cambió en mí, qué emociones manejé, cómo me relaciono con la ciudad y la comunidad, y qué le dejaré a quien viene detrás.
\end{softbox}

\textbf{Mi reflexión liberadora en cinco dimensiones.} Completa cada frase iniciada:

\small
\begin{tabularx}{\textwidth}{>{\bfseries\RaggedRight\arraybackslash}p{3.4cm}Y>{\RaggedRight\arraybackslash}p{4.6cm}}
\rowcolor{milcVino}\color{white}Dimensión & \color{white}\bfseries Mi respuesta (frame starter) & \color{white}\bfseries Algo que descubrí\\
1. Personal & Lo que más cambió en mí fue \linead{6.5cm} & \linead{4.2cm}\\
2. Emocional & La emoción más fuerte fue \linead{2.5cm} y la regulé \linead{3cm} & \linead{4.2cm}\\
3. Ciudadana & En democracia esto importa porque \linead{6cm} & \linead{4.2cm}\\
4. Local (Cartago) & En mi barrio sería útil para \linead{6.5cm} & \linead{4.2cm}\\
5. Intergeneracional & A mi abuela/abuelo le diría \linead{6.5cm} & \linead{4.2cm}\\
\end{tabularx}
\normalsize

\begin{infoband}{milcVino}
\textbf{Para la próxima sustentación.} Vuelve a esta tabla en seis meses. Verás que las respuestas cambian — no porque lo aprendido cambie, sino porque tú cambias. La evaluación liberadora es una conversación contigo a través del tiempo.
\end{infoband}


\puente{Cerramos con tres voces: la comunidad que se funda escuchando, el cambio permanente de las cosas, y la distancia que hoy significa otra cosa.}

\milcestacion{milcGradeDark}{Triángulo de pensamiento · cierre filosófico}

\begin{softbox}{milcVino}{milcGris}{Dussel · lente del nosotros}
\textit{«Aquel que es capaz de escuchar silenciosamente al Otro como otro es el que puede constituir una comunidad y no una mera sociedad totalizada.»}

{\footnotesize\itshape\color{milcNegro!70}--- Enrique Dussel · Filosofía de la liberación (1977)}

\emph{Una red de paisanos es una comunidad en el sentido exacto de Dussel:} \textbf{gente que recibe al que llega \emph{como otro}} ---distinto, con su historia, sin pedirle que se vuelva igual para ser aceptado---. \textbf{Y fíjate en lo que hace posible:} \emph{quien encuentra eso al llegar puede seguir siendo quien es mientras se adapta}, \textbf{mientras que quien llega sin nadie tiene que reinventarse solo, y muchas veces se reinventa hacia lo que cree que se espera de él.} \emph{Y hay una obligación que sale de aquí, y va dirigida a ustedes:} \textbf{el que ya está establecido decide cómo le va al que llega.} \emph{Cuando estés allá, y aparezca alguien de acá, tú eres la red.}

\textbf{Cuando alguien nuevo llegó a mi curso o a mi barrio, ¿yo fui la red o fui el que miraba?}
\end{softbox}
\linead{\linewidth}\\[1mm]
\linead{\linewidth}

\begin{softbox}{milcMostaza}{milcGris}{Estoicismo · Marco Aurelio · Meditaciones VII, 47 (c. 175 d.C.) · cuidado interior}
\textit{«Contempla las evoluciones de los astros y piensa al mismo tiempo que evolucionas con ellos; y reflexiona continuamente cómo los elementos cambian unos en otros.»}

{\footnotesize\itshape\color{milcNegro!70}--- Marco Aurelio · Meditaciones VII, 47 (c. 175 d.C.)}

\emph{A los diecisiete, irse o quedarse se siente como una decisión \textbf{definitiva}: como si uno estuviera escogiendo la vida entera en un formulario.} \textbf{Marco Aurelio insiste en lo contrario:} \emph{todo cambia, incluido uno}. \textbf{Y en este caso no es consuelo abstracto ---es un dato verificable---:} \emph{muchísima gente se fue y volvió, o se quedó y después se fue, o se fue a un lugar y terminó en otro.} \textbf{La decisión que estás tomando es más pequeña de lo que se siente,} \emph{no porque no importe, sino porque va a haber otras.} \textbf{Y hay algo que sí cambia poco:} \emph{la gente que te sostiene}. \emph{Por eso el balance de hoy se hace con nombres.}

\textbf{¿Estoy tratando esta decisión como si fuera la última que voy a tomar?}
\end{softbox}
\linead{\linewidth}\\[1mm]
\linead{\linewidth}

\begin{softbox}{milcTurquesa}{milcGris}{Floridi · lente de la infoesfera}
\textit{«Las tecnologías digitales están re-ontologizando la naturaleza misma de la infoesfera.»}

{\footnotesize\itshape\color{milcNegro!70}--- Luciano Floridi · La cuarta revolución (2014; trad.)}

\textbf{Irse cambió de naturaleza, y esta generación es la primera que lo vive así.} \emph{Antes, irse era \textbf{cortar}: cartas que tardaban semanas, una llamada cara al mes.} \textbf{Hoy uno puede irse y seguir hablando todos los días con los de acá,} \emph{y eso es una ventaja enorme ---y también una trampa---.} \textbf{La trampa tiene nombre:} \emph{se puede vivir en un lugar con la cabeza en otro}, \textbf{y así no se construye red en ninguno de los dos.} \emph{Los primeros meses en un sitio nuevo son incómodos y aburridos, y el teléfono ofrece la salida perfecta ---la anestesia del momento 6 del grado 10, aplicada al desarraigo---.} \textbf{Regla práctica:} \emph{hablar con los de allá está bien; hacerlo en vez de salir a conocer gente es lo que alarga la soledad.}

\textbf{Si me fuera, ¿cuánto de mi día seguiría estando aquí?}
\end{softbox}
\linead{\linewidth}\\[1mm]
\linead{\linewidth}

\begin{infoband}{milcNegro}
\textbf{Nota para el docente.} Las tres son citas textuales y llevan su referencia debajo. Puedes remitir a la obra en clase.
\end{infoband}

\milcestacion{milcVino}{Mi autoevaluación · marca una casilla por fila}

{\small
\setlength{\extrarowheight}{2.2pt}
\begin{tabularx}{\textwidth}{>{\RaggedRight\arraybackslash\bfseries}p{3.0cm}YYY}
\rowcolor{milcVino}\color{white}\bfseries Criterio & \color{white}\bfseries Lo logré & \color{white}\bfseries En proceso & \color{white}\bfseries Necesito apoyo\\[.6mm]
Comprensión del tema & \milcopcion{Explico las ideas con mis palabras.} & \milcopcion{Reconozco las ideas, falta precisión.} & \milcopcion{Necesito repasar lo básico.}\\[.8mm]
\rowcolor{milcGris}Calidad del balance & \milcopcion{Distingo lo que pesa de verdad —la red— de la distancia, y sé qué tendría que construir allá.} & \milcopcion{Hago las columnas, pero solo pienso en lo económico y lo académico.} & \milcopcion{Sigo pensando que irse es siempre progresar y quedarse es siempre conformarse.}\\[.8mm]
Aplicación y producto & \milcopcion{Mi producto cumple los criterios.} & \milcopcion{Está armado, con ajustes pendientes.} & \milcopcion{Aún está en construcción.}\\[.8mm]
\rowcolor{milcGris}Reflexión 5 dimensiones & \milcopcion{Respondí cada una con honestidad.} & \milcopcion{Respondí algunas con detalle.} & \milcopcion{Mis respuestas son muy generales.}\\[.8mm]
Triángulo de pensamiento & \milcopcion{Conecté las 3 voces con mi trabajo.} & \milcopcion{Conecté al menos una.} & \milcopcion{No las relaciono con mi proceso.}\\[.8mm]
\end{tabularx}}

\vspace{3mm}
\textbf{Hoy entendí que:}\\[1mm]
\linead{\linewidth}\\[1mm]
\linead{\linewidth}

\textbf{Mi compromiso para la próxima sesión es:}\\[1mm]
\linead{\linewidth}\\[1mm]
\linead{\linewidth}

\begin{softbox}{milcMostaza}{milcGris}{Cita de despedida · Marco Aurelio}
\textit{``El obstáculo en el camino se vuelve el camino. Lo que estorba al camino se vuelve camino.''}\\[1mm]
Cada error de hoy, bien observado, se convierte en pieza del camino que pisarás mañana.
\end{softbox}

\small
\begin{tabularx}{\textwidth}{>{\bfseries}p{3.6cm}Y}
\rowcolor{milcGradeDark!12}Nombre completo: & \linead{10cm}\\
Fecha: & \linead{4cm} \quad Curso/grupo: \linead{4cm}\\
Mi firma: & \linead{9cm}\\
\end{tabularx}
\normalsize

\begin{infoband}{milcGradeDark}
\textbf{Para la próxima sesión.} Dos cosas. \textbf{Primero, ten la conversación} ---con alguien que se fue y, si puedes, con alguien que se quedó pudiendo irse--- y trae \textbf{las dos versiones}. \emph{Y haz algo pequeño que cambia el cálculo: si hay alguien de tu colegio o de tu familia en el lugar al que te irías, \textbf{escríbele}. No para pedir nada ---para que exista el contacto. Eso es empezar a construir la red antes de necesitarla, que es exactamente lo que hicieron los que llegaron segundos.} \textbf{Segundo, pregunta en tu casa por la llegada:} averigua con alguien mayor \textbf{si alguien de la familia llegó de otra parte} ---de dónde, por qué, cuándo---, \textbf{quién los recibió aquí} y \textbf{qué fue lo más difícil del primer año}. \textbf{Trae:} las respuestas y lo que te contaron. \emph{Vas a encontrar la historia de este territorio entero: casi todas las familias del Norte del Valle llegaron de alguna parte, casi todas fueron recibidas por alguien, y casi nadie cuenta esa parte porque la da por hecha. Preguntar por ella es recuperar la mitad de la historia que no se cuenta.}
\end{infoband}

\end{document}
